\documentclass[10pt]{extarticle}

% ============================================================
% Pacotes
% ============================================================
\usepackage[utf8]{inputenc}
\usepackage[T1]{fontenc}
\usepackage[brazil]{babel}
\usepackage{amsmath,amssymb}
\usepackage{enumitem}
\usepackage{xcolor}
\usepackage[margin=1.5cm]{geometry}
\usepackage{titlesec}
\usepackage{fancyhdr}
\usepackage{hyperref}

% ============================================================
% Paleta de cores
% ============================================================
\definecolor{darkblue}{RGB}{0,51,102}
\definecolor{accentorange}{RGB}{230,126,34}

% ============================================================
% Estilo de títulos
% ============================================================
\titleformat{\section}
  {\normalfont\Large\bfseries\color{darkblue}}
  {\thesection}{1em}{}
\titleformat{\subsection}
  {\normalfont\large\bfseries\color{accentorange}}
  {\thesubsection}{1em}{}

% ============================================================
% Cabeçalho e rodapé
% ============================================================
\pagestyle{fancy}
\fancyhf{}
\renewcommand{\headrulewidth}{0.4pt}
\renewcommand{\footrulewidth}{0.4pt}
\fancyhead[L]{\color{darkblue}\small Problem Set 2}
\fancyhead[R]{\color{darkblue}\small PCC190 - Secure Machine Learning}
\fancyfoot[C]{\thepage}

% ============================================================
% Estilo da lista de questões
% ============================================================
\newlist{questoes}{enumerate}{1}
\setlist[questoes,1]{label=\textbf{\color{darkblue}\arabic*.},leftmargin=1.5cm,itemsep=0.8em}

\newlist{itens}{enumerate}{1}
\setlist[itens,1]{label=\alph*),leftmargin=1cm}

\begin{document}

% ============================================================
% Cabeçalho do documento
% ============================================================
\begin{center}
    {\LARGE\bfseries\color{darkblue}Problem Set 2}\\[0.3cm]
    {\large\color{accentorange} PCC190 - Secure Machine Learning}\\[0.5cm]
\end{center}


% ============================================================
% Questões
% ============================================================
\begin{questoes}

\item Define the symbol $=$.

\item Define the symbol $\triangleq$.

\item Present the formal definition of a learning algorithm given by Mitchell (1997).

\item Explain the learning framework depicted in Figure 2.1(a).

\item Define inductive bias in the context of learning algorithms. Present some examples.

\item What is empirical risk? Why can a machine learning algorithm be described as an empirical risk minimization procedure?

\item Explain the sentence: ``Minimizing the average loss on the training data is often used as a surrogate for minimizing the expected loss on unseen data.''

\item What is the stationarity assumption?

\item Define each of the symbols below: $\omega$, $\Omega$, $x$, $\mathcal{X}$, $D$, $x_i$, $\mathcal{X}_i$.

\item How does the book define data collection?

\item What is the common assumption made by all the learning algorithms studied in the book regarding the structure of the data?

\item What is the difference between explanatory variables and response variables? How do these concepts relate to the data point $x$, the label $y$, and the paired space $\mathcal{Z} = \mathcal{X} \times \mathcal{Y}$?

\item What is feature selection, and how does it differ from the measurement process $\xi$? Why does the book choose not to distinguish between the two?

\item What is a hypothesis space $\mathcal{F}$? Illustrate with an example how assumptions about the form of the hypotheses can restrict this space.

\item What is the difference between a learning model and a training procedure? Use the location estimation example (mean versus median) to explain why this distinction is useful for understanding a learner's vulnerabilities.

\item Describe the supervised learning framework. How is a learning algorithm $H^{(N)}$ formally defined, and what does it mean for such an algorithm to have a randomized element?

\item What is the difference between batch learning and online learning?

\item In a binary classification setting applied to security problems, what do the positive and negative classes typically represent? How does the anomaly detection setting studied in Chapter 6 differ from standard binary classification?

\item What is the difference between the risk $R(P_{\mathcal{Z}}, f)$ and the empirical risk $\tilde{R}_N(f)$? Why does the learner minimize the latter instead of the former?

\item What is regularization, and what problem does it address? In Equation 2.1, what do $\rho(f)$ and $\lambda$ represent?

\item What are false positives (FP) and false negatives (FN)? What are the false-positive rate (FPR) and the false-negative rate (FNR), and why is choosing a tradeoff between them an application-specific decision?

\item Why does the book restrict its analysis primarily to binary classification, and which other learning paradigms are mentioned but left out of the discussion?

\end{questoes}

\end{document}
